\documentclass[10pt,a4paper]{amsart}
\usepackage{amsmath,amssymb,amsthm,graphicx}
% For commutative diagrams we use Xy-pic
\usepackage[all]{xy}

% We use 2cell for 2-commutative diagrams.
\xyoption{2cell}
% Theorem environments.
%
\theoremstyle{plain}
\newtheorem{theorem}[subsection]{Theorem}
\newtheorem{proposition}[subsection]{Proposition}
\newtheorem{lemma}[subsection]{Lemma}

\theoremstyle{definition}
\newtheorem{definition}[subsection]{Definition}
\newtheorem{example}[subsection]{Example}
\newtheorem{exercise}[subsection]{Exercise}
\newtheorem{situation}[subsection]{Situation}

\theoremstyle{remark}
\newtheorem{remark}[subsection]{Remark}
\newtheorem{remarks}[subsection]{Remarks}

\numberwithin{equation}{subsection}

% Macros
%
\def\lim{\mathop{\mathrm{lim}}\nolimits}
\def\colim{\mathop{\mathrm{colim}}\nolimits}
\def\Spec{\mathop{\mathrm{Spec}}}
\def\Hom{\mathop{\mathrm{Hom}}\nolimits}
\def\Ext{\mathop{\mathrm{Ext}}\nolimits}
\def\SheafHom{\mathop{\mathcal{H}\!\mathit{om}}\nolimits}
\def\SheafExt{\mathop{\mathcal{E}\!\mathit{xt}}\nolimits}
\def\Sch{\mathit{Sch}}
\def\Mor{\mathop{\mathrm{Mor}}\nolimits}
\def\Ob{\mathop{\mathrm{Ob}}\nolimits}
\def\Sh{\mathop{\mathit{Sh}}\nolimits}
\def\NL{\mathop{N\!L}\nolimits}
\def\CH{\mathop{\mathrm{CH}}\nolimits}
\def\proetale{{pro\text{-}\acute{e}tale}}
\def\etale{{\acute{e}tale}}
\def\QCoh{\mathit{QCoh}}
\def\Ker{\mathop{\mathrm{Ker}}}
\def\Im{\mathop{\mathrm{Im}}}
\def\Coker{\mathop{\mathrm{Coker}}}
\def\Coim{\mathop{\mathrm{Coim}}}

% Boxtimes
%
\DeclareMathSymbol{\boxtimes}{\mathbin}{AMSa}{"02}

%
% Macros for moduli stacks/spaces
%
\def\QCohstack{\mathcal{QC}\!\mathit{oh}}
\def\Cohstack{\mathcal{C}\!\mathit{oh}}
\def\Spacesstack{\mathcal{S}\!\mathit{paces}}
\def\Quotfunctor{\mathrm{Quot}}
\def\Hilbfunctor{\mathrm{Hilb}}
\def\Curvesstack{\mathcal{C}\!\mathit{urves}}
\def\Polarizedstack{\mathcal{P}\!\mathit{olarized}}
\def\Complexesstack{\mathcal{C}\!\mathit{omplexes}}
% \Pic is the operator that assigns to X its picard group, usage \Pic(X)
% \Picardstack_{X/B} denotes the Picard stack of X over B
% \Picardfunctor_{X/B} denotes the Picard functor of X over B
\def\Pic{\mathop{\mathrm{Pic}}\nolimits}
\def\Picardstack{\mathcal{P}\!\mathit{ic}}
\def\Picardfunctor{\mathrm{Pic}}
\def\Deformationcategory{\mathcal{D}\!\mathit{ef}}
\usepackage{cleveref}
\crefname{originallemma}{original lemma}{original lemmas}
\Crefname{originallemma}{Original lemma}{Original lemmas}
\makeatletter
\newlabel{algebra-lemma-colimit-syntomic}{{07BV}{source}{}{sourcepointer-1}{}}
\newlabel{algebra-lemma-colimit-syntomic@cref}{{[originallemma][0][]07BV}{[1][0][]source}{}{}{}}
\makeatother
\crefname{originallemma}{original lemma}{original lemmas}
\Crefname{originallemma}{Original lemma}{Original lemmas}
\makeatletter
\newlabel{algebra-lemma-p-ring-map}{{0BRA}{source}{}{sourcepointer-2}{}}
\newlabel{algebra-lemma-p-ring-map@cref}{{[originallemma][0][]0BRA}{[1][0][]source}{}{}{}}
\makeatother
\crefname{originalreference}{original reference}{original references}
\Crefname{originalreference}{Original reference}{Original references}
\makeatletter
\newlabel{lemma-smooth-base-change-general}{{0F09}{source}{}{sourcepointer-3}{}}
\newlabel{lemma-smooth-base-change-general@cref}{{[originalreference][0][]0F09}{[1][0][]source}{}{}{}}
\makeatother
\crefname{originalreference}{original reference}{original references}
\Crefname{originalreference}{Original reference}{Original references}
\makeatletter
\newlabel{proposition-topological-invariance}{{03SI}{source}{}{sourcepointer-4}{}}
\newlabel{proposition-topological-invariance@cref}{{[originalreference][0][]03SI}{[1][0][]source}{}{}{}}
\makeatother
\crefname{originalreference}{original reference}{original references}
\Crefname{originalreference}{Original reference}{Original references}
\makeatletter
\newlabel{theorem-smooth-base-change}{{0EYU}{source}{}{sourcepointer-5}{}}
\newlabel{theorem-smooth-base-change@cref}{{[originalreference][0][]0EYU}{[1][0][]source}{}{}{}}
\makeatother
\title{Cohomology is invariant under extension of separably closed base fields for torsion coefficients invertible in the base}
\author{The Stacks Project Authors}
\date{stacks-project:etale-cohomology.tex:local-snapshot-20261011; GNU FDL 1.2 or later}
\begin{document}
\maketitle
\noindent\textit{Editorial note.} All delivered mathematical text is the author's own, copied from the pinned Stacks Project chapter file etale-cohomology.tex; no statement or proof step was rewritten, reordered or removed.\par\medskip
\noindent\textit{Editorial note.} The counted proof is the author's complete argument: the proof is written out in full: the reduction to algebraically closed base fields through topological invariance, the smooth-base-change step using the colimit-of-smooth-algebras description, and the passage from cohomology of quasi-compact quasi-separated objects to the isomorphism of sheaves.\par\medskip
\noindent\textit{Editorial note.} The label lemma-smooth-base-change-separably-closed names the counted statement, and the proof environment immediately following it in the source is delivered verbatim as its proof.\par\medskip
\noindent\textit{Editorial note.} The counted statement is carried by the author's own lemma environment in the chapter source.\par\medskip
\noindent\textit{Editorial note.} The statement and proof are copied from the pinned chapter file; the five cited inputs are kept as external pointers with their canonical Stacks tags.\par\medskip
\noindent\textit{Editorial note.} The preamble is the chapter's own native preamble excerpt already used for this source: the Xy-pic setup (native preamble lines 14-18) and the native theorem declarations with the mathematical macros (native preamble lines 154-217).\par\medskip
\noindent\textit{Allowed standard prerequisites (editorial).}\begin{itemize}
\item Topological invariance of etale cohomology under universal homeomorphisms. Same chapter, the proposition applying topological invariance to purely inseparable field extensions. Supplies the reduction from separably closed to algebraically closed base fields.
\item Algebraically closed fields as limits of smooth algebras. Algebra, the lemma presenting an extension of algebraically closed fields as a filtered colimit of smooth algebras. Supplies the smoothness input needed to apply the smooth base change theorem in its general form.
\item Smooth base change and its general reformulation. Same chapter, the smooth base change theorem and the general base change lemma used here. Supplies the isomorphism of cohomology along X\_K -> X under the two hypotheses on the extension.
\end{itemize}\medskip
\ifdefined\hypertarget\hypertarget{sourcepointer-1}{}\fi
\noindent\textit{Source reference.} Original lemma 07BV: sources/stacks-chow/tags.txt, source lines 9355--9355. Prerequisite pointer invoked at proof line(s) 18092; it names a standard fact of a sibling chapter that the counted argument uses, and every step of the counted proof built on it is retained verbatim in the delivered body.\par\medskip
\ifdefined\hypertarget\hypertarget{sourcepointer-2}{}\fi
\noindent\textit{Source reference.} Original lemma 0BRA: sources/stacks-chow/tags.txt, source lines 14763--14763. Prerequisite pointer invoked at proof line(s) 18084; it names a standard fact of a sibling chapter that the counted argument uses, and every step of the counted proof built on it is retained verbatim in the delivered body.\par\medskip
\ifdefined\hypertarget\hypertarget{sourcepointer-3}{}\fi
\noindent\textit{Source reference.} Original reference 0F09: sources/stacks-etale-cohomology/pinned-original-source.tex, source lines 17975--17975. Prerequisite pointer invoked at proof line(s) 18095; it names a standard fact of the same chapter that the counted argument uses, and every step of the counted proof built on it is retained verbatim in the delivered body.\par\medskip
\ifdefined\hypertarget\hypertarget{sourcepointer-4}{}\fi
\noindent\textit{Source reference.} Original reference 03SI: sources/stacks-etale-cohomology/pinned-original-source.tex, source lines 6548--6548. Prerequisite pointer invoked at proof line(s) 18088; it names a standard fact of the same chapter that the counted argument uses, and every step of the counted proof built on it is retained verbatim in the delivered body.\par\medskip
\ifdefined\hypertarget\hypertarget{sourcepointer-5}{}\fi
\noindent\textit{Source reference.} Original reference 0EYU: sources/stacks-etale-cohomology/pinned-original-source.tex, source lines 17446--17446. Prerequisite pointer invoked at proof line(s) 18094; it names a standard fact of the same chapter that the counted argument uses, and every step of the counted proof built on it is retained verbatim in the delivered body.\par\medskip
\begin{lemma}
\label{lemma-smooth-base-change-separably-closed}
Let $K/k$ be an extension of separably closed fields. Let $X$
be a quasi-compact and quasi-separated scheme over $k$.
Let $E \in D^+(X_\etale)$ have cohomology sheaves whose stalks
are torsion of orders invertible in $k$. Then
\begin{enumerate}
\item the maps $H^q_\etale(X, E) \to H^q_\etale(X_K, E|_{X_K})$
are isomorphisms, and
\item $E \to R(X_K \to X)_*E|_{X_K}$ is an isomorphism.
\end{enumerate}
\end{lemma}

\begin{proof}
Proof of (1). First let $\overline{k}$ and $\overline{K}$
be the algebraic closures
of $k$ and $K$. The morphisms $\Spec(\overline{k}) \to \Spec(k)$ and
$\Spec(\overline{K}) \to \Spec(K)$ are universal homeomorphisms
as $\overline{k}/k$ and $\overline{K}/K$ are purely inseparable
(see Algebra, Lemma \ref{algebra-lemma-p-ring-map}).
Thus $H^q_\etale(X, \mathcal{F}) =
H^q_\etale(X_{\overline{k}}, \mathcal{F}_{X_{\overline{k}}})$ by
the topological invariance of \'etale cohomology, see
Proposition \ref{proposition-topological-invariance}.
Similarly for $X_K$ and $X_{\overline{K}}$.
Thus we may assume $k$ and $K$ are algebraically closed.
In this case $K$ is a limit of smooth $k$-algebras, see
Algebra, Lemma \ref{algebra-lemma-colimit-syntomic}.
We conclude our lemma is a special case of
Theorem \ref{theorem-smooth-base-change} as reformulated in
Lemma \ref{lemma-smooth-base-change-general}.

\medskip\noindent
Proof of (2). For any quasi-compact and quasi-separated $U$ in $X_\etale$
the above shows that the restriction of the map
$E \to R(X_K \to X)_*E|_{X_K}$ determines an isomorphism on cohomology.
Since every object of $X_\etale$ has an \'etale covering by such $U$
this proves the desired statement.
\end{proof}

\end{document}
